\chapter{Industrialisation in India}

\section{Industrialisation and De-industrialisation}

\begin{KeyConcept}
\begin{itemize}
\item \textbf{Industrialisation} = producing on a large/mass scale, with an organised production process and labour drawn from different domains (not just household/kinship). It is studied by economics (GDP, employment), history (how it happened) and sociology (through social categories --- caste, gender, discrimination, exploitation).
\item \textbf{De-industrialisation} (propounded by \textbf{William Digby and R. C. Dutt}) is based on three propositions: (1) the \textbf{decline of handicraft industry} continued well into the 20th century; (2) it was \textbf{not compensated} by a sufficient rise in modern industries; (3) in consequence, \textbf{Indian society became more and more agricultural} (the displaced handicraft workers became agricultural labourers).
\item At \textbf{1800} India had about \textbf{15--20\% of its workforce in manufacturing} --- a very large manufacturing population.
\item The idea of de-industrialisation has two roots: the \textbf{nationalist tradition} (R. C. Dutt, \textbf{Dadabhai Naoroji}, and later Nehru) and the \textbf{Marxist theory of imperialism} (the colonial powers first captured the economy, then the political and social domains). It is a precursor of \textbf{dependency theory} (A. G. Frank).
\end{itemize}
\end{KeyConcept}

\section{The De-industrialisation Debate: Thorner vs Krishnamurthy}

\begin{KeyConcept}
\begin{itemize}
\item \textbf{Daniel Thorner} (a Marxist) --- there was \textbf{no de-industrialisation}. Using census data (1881--1931), he found the \textbf{male industrial working force roughly constant (or increasing)}, so the census data do not support either absolute or relative de-industrialisation. If a shift from industry to agriculture happened, it was \textbf{between 1815 and 1880} --- due to the rise of \textbf{plantation agriculture} (especially in the North-East) and forced commercial crops (indigo), which drew in people who were \emph{not} in manufacturing.
\item \textbf{J. Krishnamurthy} --- critiqued Thorner: industrialisation means the \textbf{rise in the share of manufacturing output per capita} (historically linked to a rise in the share of manufacturing in the workforce, but that is \textbf{neither a necessary nor a sufficient condition}). A \textbf{decline in manufacturing output per capita = de-industrialisation}, and a declining workforce is compatible with rising output under \textbf{capital-intensive techniques}. So a constant workforce does \emph{not} disprove de-industrialisation.
\end{itemize}
\end{KeyConcept}

\begin{QuickRevision}
\begin{itemize}
\item Industrialisation = mass-scale, organised production, labour from diverse domains.
\item De-industrialisation (Digby \& R. C. Dutt): decline of handicrafts (uncompensated) $\rightarrow$ society more agricultural; 1800 = 15--20\% in manufacturing; roots = nationalism + Marxist imperialism.
\item Thorner: no de-industrialisation (census 1881--1931, workforce constant); shift 1815--80 (plantations, indigo).
\item Krishnamurthy: industrialisation = output per capita (not workforce); capital-intensive techniques; output decline = de-industrialisation.
\end{itemize}
\end{QuickRevision}

\section{Tirthankar Roy's Synthesis: Traditional vs Modern Industry}

\begin{KeyConcept}
\begin{itemize}
\item \textbf{Tirthankar Roy} summarised the debate --- he also holds there was \textbf{no de-industrialisation, but industrialisation}; the real difference is between \textbf{traditional and modern industry}.
\item \textbf{Traditional (pre-British) industries}: \textbf{spinning and weaving, leather goods, metal works, carpets and rugs} --- they used \textbf{no machinery}, were \textbf{not organised in large-scale factories}, used \textbf{family/household labour}, and were \textbf{not regulated by law}.
\item \textbf{Modern industry} (a product of the \textbf{Industrial Revolution}): uses machinery, large-scale factories, is regulated by law, and recruits from the open market.
\item The \textbf{de-industrialisation theory} (net effect of British policy was negative) rests on four propositions: (1) traditional industry declined; (2) it declined because \textbf{hand tools lost to machinery} (obsolete technology); (3) this market battle was \textbf{forced by Britain's commitment to free trade} (no subsidy/protection for handicrafts); (4) Britain did create some modern industries (textile mills) but they \textbf{did not compensate} for the destruction (Manchester competed with, and undercut, the handicrafts whose goods were actually of \emph{better} quality but produced more slowly and costlier).
\end{itemize}
\end{KeyConcept}

\section{Why and How Factories Came to India}

\begin{KeyConcept}
\begin{itemize}
\item Factories were established in India \textbf{after the 1850s}, mainly to \textbf{export manufactured goods to Britain} --- first the \textbf{textile mills of Bombay} and the \textbf{jute mills of Kolkata}, then Madras. (Plantation agriculture --- tea in Assam --- began even earlier, in \textbf{1839}, for British consumption.)
\item \textbf{Why labour was so cheap}: (1) the \textbf{destruction of handicrafts} threw skilled workers into the market; (2) \textbf{peasants forced into cash crops} (indigo) got very low wages (in textiles and plantations the wage was about \textbf{one-third}); (3) \textbf{labour was not regulated by law}, and \textbf{child labour was rampant}.
\item The British also developed \textbf{hill stations} --- the seat of power is always at the top (Raisina Hill), and controlling the hill stations gave legitimacy over the hill people --- built by local labour at a very high death rate.
\end{itemize}
\end{KeyConcept}

\begin{KeyConcept}
\begin{itemize}
\item \textbf{World War I} --- Britain needed goods it could no longer produce at home, so it \textbf{liberalised the rules} and allowed Indians to establish factories (the \textbf{steel plant at Jamshedpur}, ordnance factories). \textbf{D. R. Gadgil} says industrialisation really began during World War I.
\item \textbf{D. R. Gadgil's point}: World War I led to the \textbf{expansion of the working class}, and --- because Indians could now be \textbf{promoted from worker to supervisor} (the 'white-collar' jobs, so called because supervisors wore white shirts, earlier reserved for the British) --- it created \textbf{multiple classes}, and even \textbf{attracted the upper castes} into industry.
\item The \textbf{Indian Railways} (1860s) also pushed industrialisation --- first linking the textile and jute mills, then (in the early 20th century) expanding and aiding indigenous factories.
\end{itemize}
\end{KeyConcept}

\section{The Colonial Working Class: Morris D. Morris and Caste in the Mills}

\begin{KeyConcept}
\begin{itemize}
\item \textbf{Morris D. Morris} studied the working class of the Bombay textile mills (colonial period). The \textbf{Factory Acts of 1881 and 1891} (against child labour, on working hours) were passed under pressure from \textbf{Manchester} (which could not compete with Indian mills), but were \textbf{ineffective} --- the offices were controlled by British officials, there was corruption, and \textbf{child labour remained rampant}.
\item Morris's findings: workers of \textbf{different castes formed clusters according to their status group} (same caste = same status, living together), but \textbf{recruitment was not based on caste}. However, \textbf{Dalits were systematically excluded from the spinning section} of the textile mills --- because spinning required using a spit (saliva), and the \textbf{spinning section had the higher wages} --- so caste norms kept Dalits out of the most lucrative section.
\end{itemize}
\end{KeyConcept}

\begin{QuickRevision}
\begin{itemize}
\item Tirthankar Roy: traditional (no machinery, household labour, unregulated) vs modern (machinery, large factory, regulated, open-market) industry.
\item De-industrialisation's 4 propositions: decline, obsolete tech (hand tools lost), free trade, modern industries didn't compensate.
\item Factories after 1850s (Bombay textiles, Kolkata jute, Madras); plantation (Assam tea 1839); cheap labour (destruction of handicrafts, forced cash crops, no regulation, child labour).
\item Hill stations = power at the top (Raisina Hill).
\item WWI: Gadgil --- liberalisation $\rightarrow$ factories (Jamshedpur); worker $\rightarrow$ supervisor (white collar) $\rightarrow$ multiple classes; upper castes entered; Railways (1860s).
\item Morris D. Morris: Factory Acts 1881/1891 ineffective; caste clusters; Dalits excluded from the spinning section (spit, higher wages).
\end{itemize}
\end{QuickRevision}

